% Vietnamese translation for OI Public Library
% Source-File: IOI中国国家候选队论文/2004/楼天城.ppt
\documentclass[11pt]{article}
\usepackage{oipl-vn}

\newcommand{\Oh}{\mathcal{O}}

\title{Tìm kiếm bộ phận kết hợp thuật toán hiệu quả\\{\large Ghi chú theo slide}}
\author{Lou Tiancheng}
\date{IOI 2004}

\begin{document}
\maketitle

\origtitle{Partial search plus efficient algorithms}
\sourcepath{IOI 2004 / Lou Tiancheng.ppt}

\section{Động cơ}

Các slide đặt vấn đề bằng những bài có không gian tìm kiếm cỡ \(\Oh(n!)\).
Nếu toàn bộ cấu hình đều phải chọn bằng DFS thì không thể chạy kịp, nhưng
nhiều phần của cấu hình lại có thể giải bằng ghép cặp, quy hoạch động hoặc
kiểm tra hình học sau khi cố định một số biến nhỏ.

\section{Mô hình tìm kiếm bộ phận}

Thay vì tìm kiếm mọi đối tượng, ta chỉ tìm trên phần gây ràng buộc khó nhất.
Phần còn lại được giao cho một thuật toán hiệu quả. Bản chất là tách biến:
nhóm biến nhỏ được liệt kê, nhóm biến lớn được tối ưu bằng cấu trúc đã biết.

Các slide nhắc ví dụ với bộ \((a,b,c,d)\), các dòng số như
\[
  (1,3,5,6),\quad (2,2,5,3),\quad (3,1,4,1),
\]
và bài \emph{Triangle Construction}. Với bài hình học này, sau khi cố định
một số điểm hoặc đoạn then chốt, phần kiểm tra còn lại có thể dùng quan hệ
khoảng cách và ràng buộc bán kính thay vì thử mọi cách ghép.

\section{Ví dụ NOI 2003}

Một mốc khác là bài \emph{Phá trận liên hoàn}: các slide nhắc điều kiện
\[
  (u_i-x_j)^2+(v_i-y_j)^2\le R^2
\]
và bảng \(C[S][T][I]\) cho biết mục tiêu \(I\) có thể phủ bởi đoạn
\([S,T]\) hay không. Sau khi tiền xử lý bảng này trong khoảng \(\Oh(n^3)\),
việc tìm kiếm chỉ còn xét các khoảng ứng viên \([S_i,T_i]\) thỏa quan hệ
ghép nối.

\section{Kết luận}

Phương pháp này hữu ích khi một bài ``trông như tìm kiếm'' nhưng có những
khối con rất đều. Tách đúng khối con sẽ giảm nhanh số nhánh, đồng thời giữ
được chứng minh rõ ràng cho phần được xử lý bằng thuật toán chuyên biệt.

\end{document}
